% Generated draft: requires mathematical and source review.
\documentclass[11pt]{article}
\usepackage[margin=25mm]{geometry}
\usepackage{fontspec}
\setmainfont{Latin Modern Roman}
\usepackage{amsmath,amssymb,mathtools,unicode-math}
\setmathfont{Latin Modern Math}
\usepackage{enumitem,needspace,xurl,hyperref}
\setlength{\parindent}{0pt}
\setlength{\parskip}{5pt}
\setlength{\emergencystretch}{3em}
\newcommand{\N}{\mathbb{N}}
\newcommand{\Z}{\mathbb{Z}}
\newcommand{\Q}{\mathbb{Q}}
\newcommand{\R}{\mathbb{R}}
\newcommand{\C}{\mathbb{C}}
\newcommand{\F}{\mathbb{F}}
\newcommand{\A}{\mathbb{A}}
\newcommand{\PP}{\mathbb{P}}
\newcommand{\E}{\mathbb{E}}
\newcommand{\eps}{\varepsilon}
\newcommand{\dd}{\,\mathrm{d}}
\begin{document}
\section*{UnsolvedMath ID: 20002690}
\textbf{Source problem code:} catalogue entry
\subsection{Definitions and mathematical statement}
\textbf{Source question.} Define an \(R\)-transform for an unbounded pair \((x_1,x_2)\), since a moment power series may not exist or determine its law. The source suggests the freely infinitely divisible case as an easier subproblem. It does not specify self-adjointness, affiliation, joint-law equivalence, analytic domains or a proposed multivariable transform. Scalar and genuinely noncommutative pair formulations must be distinguished.

\textbf{Reported scalar analytic theory.} For a probability measure \(\mu\) on \(\R\), define
\[
 G_\mu(z)=\int_\R\frac{\mathrm d\mu(t)}{z-t},\qquad
 F_\mu(z)=G_\mu(z)^{-1}.
\]
The aid attributes to Bercovici--Voiculescu an inverse branch on suitable truncated Stolz cones \(\{z:\operatorname{Im}z>\beta,\ |\operatorname{Re}z|<\alpha\operatorname{Im}z\}\), giving
\[
 \phi_\mu(z)=F_\mu^{-1}(z)-z,\qquad R_\mu(w)=\phi_\mu(1/w)
\]
on the corresponding domains. No finite moments or global Taylor series are required. Freely infinitely divisible laws are reported to admit an upper-half-plane extension with the free L\'evy--Khintchine representation
\[
 \phi_\mu(z)=\gamma+\int_\R\frac{1+tz}{z-t}\,\mathrm d\sigma(t),
 \qquad\gamma\in\R,\quad\sigma\text{ finite positive}.
\]
This is an extension of \(\phi\) on the upper half-plane, not an assertion that every \(R\)-transform is globally analytic there or that arbitrary unbounded laws have such a representation.

\textbf{Reported matrix-valued pair reduction.} Let \((M,\tau)\) be finite tracial, and let \(X_1,X_2\) be affiliated operators. Interpret the following block in the algebra of measurable affiliated matrix operators, with its self-adjoint closed realization and domain issues checked:
\[
 L_X=\begin{pmatrix}
 0&X_1&X_2\\X_1^*&0&0\\X_2^*&0&0
 \end{pmatrix}.
\]
Equivalently one may use the self-adjoint star block of the closed row operator. This treats even non-self-adjoint coordinates and does not silently assume that an arbitrary formal unbounded matrix sum is self-adjoint. Regard \(L_X\) as an \(M_3(\C)\)-valued variable under \(E=\operatorname{id}\otimes\tau\). Its fully matricial resolvent data are
\[
 G_X^{(k)}(b)=(\operatorname{id}_{M_{3k}}\otimes\tau)
 \left[(b-I_k\otimes L_X)^{-1}\right],
 \qquad \operatorname{Im}b>0,\quad k\geq1.
\]
The aid invokes Williams's finite-dimensional operator-valued unbounded analytic theory to obtain the corresponding \(R\)-transform on its valid inversion domains. It reports recovery of the ordered pair through the fixed matrix corners, up to the appropriate trace-preserving spatial isomorphism of the generated von Neumann algebras and their affiliated operators. Full analytic recovery/inverse-domain hypotheses are not supplied here and require verification; a scalar restriction of the matrix argument is not substituted for the fully matricial data.

For freely independent pairs \(X,Y\), meaning their generated algebras are free, their block variables are free with amalgamation over \(M_3(\C)\). The aid reports, on common valid analytic domains and with affiliated sums interpreted by closure,
\[
 L_{X+Y}=L_X+L_Y,\qquad
 R_{L_{X+Y}}=R_{L_X}+R_{L_Y}.
\]
This is the proposed componentwise-addition linearization. It does not supply a universal multivariable L\'evy--Khintchine classification merely from the scalar one.

\textbf{Reported obstruction to directional scalar data.} For Hermitian matrices \(A,B\), every real pencil \(sA+tB\) has the same scalar spectral law as \(sA^T+tB^T\), since the latter is its transpose. The aid reports an explicit \(3\times3\) Hermitian example for which some ordered sixth moment differs between \((A,B)\) and \((A^T,B^T)\). The matrix entries and the particular word are not included in the supplied summary, so no explicit example is invented. This illustrates why scalar directional laws can miss noncommutative ordering information. The analytic matrix-valued route is a reported reduction, and its full unbounded-domain and information-recovery claims remain subject to primary-source review.
\subsection{Short English statement}
The scalar unbounded transform uses analytic inverse Cauchy transforms rather than moments. The aid proposes encoding an affiliated pair in a self-adjoint matrix block and using fully matricial operator-valued transforms to retain ordering and linearize free addition. Scalar directional transforms alone can lose ordered-moment information.
\subsection{Sources}
\begin{itemize}
\item Local UnsolvedMath registry, permanent ID \(20002690\), catalogue entry, \texttt{aim-probability-notes.json}, \texttt{notes[131]}, \emph{Free analysis}, subsection 0.10 ``Infinite Divisibility'', second question, printed page 7, source item 30: \url{https://aimath.org/WWN/freeanalysis/freeanalysis.pdf}.
\item Supplied research attributions: Bercovici--Voiculescu scalar unbounded analytic transforms and free L\'evy--Khintchine theory; Williams's finite-dimensional operator-valued unbounded analytic theory. No titles, dates or full theorem statements are supplied.
\item Supplied research summary on the star linearization, pair recovery/additivity and transpose-matrix obstruction. Attempt 1: partial\_result, partially\_solved, rigor ``proved'', low novelty confidence. The machine-generated aid is not peer reviewed and requires independent verification before citation. The supplied verification reports visual PDF inspection; this draft has not visited it.
\end{itemize}
\end{document}
